\todo{Lecture 21}
\begin{theorem}
Un ensemble $E\subset \mathbf{R}^{n}$ est mesurable si et seulement si $\partial E$ est negligeable.
\end{theorem}
\begin{proof}
Soit $R$ un pave qui contient $E$. 
\begin{itemize}
\item Direction $\Rightarrow$ : Supposons $E$ mesurable, alors pour tout $\varepsilon>0$, il existe $P_{\varepsilon}$ telle que 
$$
\overline{S}(\mathbb{1}_{E},P_{\varepsilon})-\underline{S}(\mathbb{1}_{E},P_{\varepsilon})<\varepsilon
$$
si et seulement si
$$
\sum _{Q\cap E\neq \varnothing, Q\cap E^{\complement}\neq\varnothing} \mathrm{Vol}(Q)<\varepsilon.
$$
On a 
$$
\partial E\subset \bigcup _{Q\subset \mathcal{E}}Q\cup \underbrace{ \bigcup _{Q\subset P_{\varepsilon}}\partial Q }_{ \text{negligeable} }
$$
ou
$$
\mathcal{E}=\{ Q:Q\cap E\neq \varnothing, Q\cap E^{\complement}\neq \varnothing \}.
$$
Soit 
$$
\mathcal{E}'=\{ Q\subset E:\partial Q,Q\subset P_{\varepsilon} \}
$$
donc
$$
\sum _{\tilde{Q}\in \mathcal{E}'}\mathrm{Vol}(\tilde{Q})<\varepsilon
$$
et comme $\partial E\subset \bigcup _{\tilde{Q}\in \mathcal{E}'}\tilde{Q}$ alors $\partial E$ est negligeable par \ref{lem811}.
\item Direction $\Leftarrow$ : Supposons $\partial E$ negligeable (c-a-d $\mathbb{1}_{\partial E}$ est integrable et $\int \mathbb{1}_{\partial E}=0$) , alors pour tout $\varepsilon>0$, il existe $P_{\varepsilon}$ telle que
$$
\overline{S}(\mathbb{1}_{\partial E},P_{\varepsilon})-\underline{S}(\mathbb{1}_{\partial E},P_{\varepsilon})<\varepsilon
$$
et donc
$$
\sum _{Q\in P_{\varepsilon},Q\cap \partial E\neq \varnothing} \mathrm{Vol}(Q)<\varepsilon.
$$
On a
$$
\overline{S}(\mathbb{1}_{E},P_{\varepsilon})-\underline{S}(\mathbb{1}_{E},P_{\varepsilon})=\sum _{\underset{Q\cap E^{\complement}\neq \varnothing}{Q\cap E\neq \varnothing }}\mathrm{Vol}(Q)\le \sum _{Q\cap \partial E\neq \varnothing}\mathrm{Vol}(Q)<\varepsilon
$$
et donc $\mathbb{1}_{E}$ est integrable et $E$ est mesurable. 
\end{itemize}
\end{proof}
\begin{remark}
Soient $E,F\subset \mathbf{R}^{n}$ mesurables. 
\begin{itemize}
\item $E\cap F$ est mesurable.
\begin{proof}
On a que $\mathbb{1}_{E\cap F}=\mathbb{1}_{E}\mathbb{1}_{F}\in \mathcal{R}(\mathbf{R}^{n})$.
\end{proof}
\item $E\cup F$ est mesurable.
\begin{proof}
On a que $\mathbb{1}_{E\cup F}=\mathbb{1}_{E}+\mathbb{1}_{F}-\mathbb{1}_{E\cap F}\in \mathcal{R}(\mathbf{R}^{n})$.
\end{proof}
\item $E\setminus F$ est mesurable.
\item $\mathring{E}$ est mesurable.
\item $\overline{E}$ est mesurable.
\end{itemize}
\end{remark}
\begin{notation}
On note $\mathcal{J}(\mathbf{R}^{n})$ l'ensemble des sous-ensembles mesurables et compacts de $\mathbf{R}^{n}$.
\end{notation}
\section{Caractérisation des fonctions intégrables}
\begin{theorem}
Soit $R\subset \mathbf{R}^{n}$ un pave et $f:R\to \mathbf{R}$ bornée. Si l'ensemble des points de discontinuité de $f$ est negligeable, alors $f$ est integrable.
\end{theorem}
\begin{proof}
Soit $N$ l'ensemble des points de discontinuité de $f$ et $M:=\sup_{\boldsymbol{x}\in R}|f(\boldsymbol{x})|$. Puisque $N$ est negligeable, on sait que 
$$
\forall\varepsilon>0\,\exists \tilde{P}_{\varepsilon}:\overline{S}(\mathbb{1}_{N},\tilde{P}_{\varepsilon})-\underbrace{ \underline{S}(\mathbb{1}_{N},\tilde{P}_{\varepsilon}) }_{ =0 }<\frac{\varepsilon}{1+4M},
$$
ou 
$$
\overline{S}(\mathbb{1}_{N},\tilde{P}_{\varepsilon})=\sum _{Q\cap N\neq \varnothing}\mathrm{Vol}(Q)<\frac{\varepsilon}{1+4M}.
$$

Soit 
$$
K=\bigcup _{\underset{Q\cap N=\varnothing}{Q\in P_{\varepsilon}}}Q.
$$
$K$ est compact, donc $f$ est uniformément continue sur $K$ :
$$
\exists\delta _{\varepsilon}>0\,\forall \boldsymbol{x},\boldsymbol{y}\in K:\lVert \boldsymbol{x}-\boldsymbol{y} \rVert <\delta _{\varepsilon}\Rightarrow |f(\boldsymbol{x})-f(\boldsymbol{y})|< \frac{\varepsilon}{1+2\mathrm{Vol}(R)}.
$$
Soit $P_{\varepsilon}$ un raffinement de $\tilde{P}_{\varepsilon}$, alors
$$
\mathrm{diam}(Q)=\sup_{\boldsymbol{x},\boldsymbol{y}\in Q}\lVert \boldsymbol{x}-\boldsymbol{y} \rVert \le \delta _{\varepsilon}.
$$
On a 
\begin{align*}
\overline{S}(f,P_{\varepsilon})-\underline{S}(f,P_{\varepsilon}) & =\sum _{Q\in P_{\varepsilon}}(\sup_{Q}f-\inf_{Q}f )\mathrm{Vol}(Q) \\
 & =\sum _{Q\subset K}\underbrace{ (\sup_{Q}f-\inf_{Q}f) }_{ < \frac{\varepsilon}{1+2\mathrm{Vol}(R)} }\mathrm{Vol}(Q)+\sum _{Q\cap N\neq \varnothing}\underbrace{ (\sup_{Q}f-\inf_{Q}f) }_{ \le 2M }\mathrm{Vol}(Q) \\
 & \le \frac{\varepsilon}{1+2\mathrm{Vol}(R)}\mathrm{Vol}(R)+2M \frac{\varepsilon}{1+4M}<\varepsilon
\end{align*}
donc $f$ est integrable.
\end{proof}
\begin{corollary}
Soit $E\subset \mathbf{R}^{n}$ negligeable et $f:E\to \mathbf{R}$ bornée. Alors $f\in \mathcal{R}(E)$ et $\int f=0$. 
\end{corollary}
\begin{proof}
Soit $R$ un pave contenant $E$ et $\tilde{f}$ le prolongement par zero. Alors l'ensemble de discontinuités de $\tilde{f}$ est negligeable et $\tilde{f}\in \mathcal{R}(R)$, donc $f\in \mathcal{R}(E)$ et $\int f=0$.
\end{proof}
\begin{theorem}
Soit $E\subset \mathbf{R}^{n}$ mesurable. Si $f:E\to \mathbf{R}$ est bornée et continue sur ${\mathring{ E }}$, alors $f$ est integrable.
\end{theorem}
\paragraph{Propriétés de l'integrale de Riemann}
\begin{enumerate}[5]
\item Soit $E\subset \mathbf{R}^{n}$ et $f\in \mathcal{R}(E)$. Alors 
$$
(\inf_{E}f)\mathrm{Vol}(E)\le \int _{E}f(\boldsymbol{x})\,d\boldsymbol{x}\le (\sup_{\boldsymbol{x}\in E}f(\boldsymbol{x}))\mathrm{Vol}(E).
$$
\item Si $E\in \mathcal{J}(\mathbf{R}^{n})$ (mesurable et compact) est connexe par arcs et $f:E\to \mathbf{R}$ est continue, alors il existe $\boldsymbol{x}_{0}\in E$ tel que
$$
\frac{1}{\mathrm{Vol}(E)}\int _{E}f(\boldsymbol{x})\,d\boldsymbol{x}=f(\boldsymbol{x}_{0}).
$$
\item Soit $E\subset \mathbf{R}^{n}$ mesurable et $f\in \mathcal{R}(E)$. Si $E=E_{1}\cup E_{2}$ tels que $E_{1}$, $E_{2}$ sont mesurables et $E_{1}\cap E_{2}$ est negligeable, alors $f|_{E_{1}}\in \mathcal{R}(E_{1})$, $f|_{E_{2}}\in \mathcal{R}(E_{2})$ et $$\int _{E}f(\boldsymbol{x})\,d\boldsymbol{x}=\int _{E_{1}}f(\boldsymbol{x})\,d\boldsymbol{x}+\int _{E_{2}}f(\boldsymbol{x})\,d\boldsymbol{x}.$$
\end{enumerate}
\begin{remark}
Supposons $E$ ferme, alors ${\mathring{ E }\cup \partial E}$ qui est l'union de deux ensemble mesurables. Alors 
$$
\int _{E}f(\boldsymbol{x})\,d\boldsymbol{x}=\int _{{\mathring{ E }}}f(\boldsymbol{x}) \, d\boldsymbol{x}+\int _{\partial E}f(\boldsymbol{x}) \, d\boldsymbol{x}  
$$
ou $\int _{\partial E}f(\boldsymbol{x}) \, d\boldsymbol{x}=0$ car $\partial E$ est negligeable.
\end{remark}
\section{Généralisation de la formule des intégrales itérées}
Soit $R=[a_{1},b_{1}]\times[a_{2},b_{2}]$ et $f:R\to \mathbf{R}$ continue. Alors
\begin{align*}
\int _{R}f(x,y)\,dx\,dy & =\int_{a_{2}}^{b_{2}} \left( \int_{a_{1}}^{b_{1}} f(x,y) \, dx  \right) \, dy \\
 &  =\int_{a_{1}}^{b_{1}} \left( \int_{a_{2}}^{b_{2}} f(x,y) \, dy  \right) \, dx .
\end{align*}
\begin{notation}
$$
\boldsymbol{z}=(\underbrace{ z_{1},\dots,z_{n} }_{ \boldsymbol{x}\in \mathbf{R}^{n} },\underbrace{ z_{n+1} }_{ y\in \mathbf{R} })\in \mathbf{R}^{n+1}
$$
\end{notation}
\begin{definition}
On dit que $E\subset \mathbf{R}^{n+1}$ est un domaine simple par rapport a la variable $y$ s'il existe un compact $K\in \mathcal{J}(\mathbf{R}^{n})$ et deux fonctions continues $g,h:K\to \mathbf{R}$ tels que $g(\boldsymbol{x})\le h(\boldsymbol{x})$ pour tout $\boldsymbol{x}\in K$ et 
$$
E=\{ (\boldsymbol{x},y)\in \mathbf{R}^{n+1}:\boldsymbol{x}\in K, g(\boldsymbol{x})\le y \le h(\boldsymbol{x}) \}.
$$
\end{definition}
\begin{theorem}
Un domaine simple $E=\{ (\boldsymbol{x},y)\in \mathbf{R}^{n+1}:\boldsymbol{x}\in K,\, g(\boldsymbol{x})\le y \le h(\boldsymbol{x}) \}$ avec $K\in \mathcal{J}(\mathbf{R}^{n})$ et $g,h\in C^{0}(K)$ est mesurable et 
$$
\mathrm{Vol}(E)=\int _{K}(h(\boldsymbol{x})-g(\boldsymbol{x}))\,d\boldsymbol{x}.
$$
\end{theorem}
\begin{theorem}
Soit $E=\{ (\boldsymbol{x},y)\in \mathbf{R}^{n+1}:\boldsymbol{x}\in K,\,g(\boldsymbol{x})\le y \le h(\boldsymbol{x}) \}$ un domaine simple et $f:E\to \mathbf{R}$ continue. Alors $f\in \mathcal{R}(E)$ et 
$$
\int _{E}f(\boldsymbol{x},y)\,d\boldsymbol{x}\,dy=\int _{K}\left( \int_{g(\boldsymbol{x})}^{h(\boldsymbol{x})} f(\boldsymbol{x},y) \, dy  \right) \, d\boldsymbol{x} 
$$
\end{theorem}
\begin{example}
Soit $E=\{ (x,y)\in \mathbf{R}^{2}:x^{2}+y^{2} \le 1 \}$. On souhaite calculer 
$$
\mathrm{Vol}(E)=\int _{E}\mathbb{1}_{E}(x,y)\,dx\,dy=\pi
$$
et
$$
I=\int _{E}(x^{2}+y^{2})\,dx\,dy.
$$
On a 
$$
E=\{ (x,y)\in \mathbf{R}^{2}:-1\le x\le 1, \, -\sqrt{ 1-x^{2} }\le y \le \sqrt{ 1-x^{2} } \},
$$
donc on applique le théorème
\begin{align*}
\mathrm{Vol}(E) & =\int _{E}\mathbb{1}_{E}(x,y)\,dx\,dy \\
 & =\int_{-1}^{1} \left( \int_{-\sqrt{ 1-x^{2} }}^{\sqrt{ 1-x^{2} }} 1 \, dy  \right) \, dx \\
  & =\int_{-1}^{1} 2\sqrt{ 1-x^{2} } \, dx =\dots=\pi
\end{align*}
et
\begin{align*}
I & =\int _{E}(x^{2}+y^{2})\,dx\,dy \\
 & =\int_{-1}^{1} \left( \int_{-\sqrt{ 1-x^{2} }}^{\sqrt{ 1-x^{2} }} x^{2}+y^{2} \, dy  \right) \, dx \\
 & =\int_{-1}^{1} 2x^{2}\sqrt{ 1-x^{2} }+ \frac{2(1-x^{2})^{3/2}}{3}  \, dx \\
 & =\dots  
\end{align*}
\end{example}